\begin{abstract}
Context resets can remove stale reasoning from large language model (LLM) agents, but can also discard useful investigation state. We present MORPH-FL, a Python fault-localization framework that separates observations, hypotheses and dialogue in four linked runtime models and supports evidence-preserving resets. We compare adaptive reset, continuous context, fixed-interval reset and an observation-only view on 50 issues from 12 Python repositories, using Qwen2.5-Coder-7B and DeepSeek-Coder-6.7B. The study contains all 400 model--workflow investigations and 100 retrieval-control records under a 32,000-token ceiling and one seed. Adaptive, continuous and fixed-reset workflows have identical per-issue file Top-1 hit indicators: 42\% with Qwen and 40\% with DeepSeek. Observation-only reaches 50\% and 40\%, respectively. Equal file accuracy conceals different investigation output: DeepSeek retains grounded hypotheses in 6/50 adaptive runs versus 47/50 observation-only runs. Adaptive resets occur in 8/50 Qwen and 29/50 DeepSeek runs without additional Top-1 hits. Observation-only takes 140.3 versus 93.6 seconds with Qwen and 111.5 versus 55.0 seconds with DeepSeek. These results show why context-policy evaluation needs output availability, reset exposure and resource use alongside rankings. Patch-derived targets, bundled configuration changes and a known prompt-growth defect constrain causal interpretation; the study does not establish that adaptive resetting improves localization.
\end{abstract}

